\section{Research Agenda}\label{sec:agenda}

The gaps of \cref{sec:synthesis} suggest fourteen projects that a networking or multimedia laboratory
could start with existing hardware and open software. We order them by dependency (\cref{fig:roadmap}).
Projects that produce shared artefacts come first, because the others need them.

\begin{figure}[t]
\centering
\begin{adjustbox}{max width=\linewidth}\begin{tikzpicture}[font=\sffamily\scriptsize,
  pj/.style={draw=black!55,rounded corners=2pt,fill=white,align=left,text width=4.45cm,inner sep=2.2pt,minimum height=6.4mm},
  ph/.style={font=\sffamily\footnotesize\bfseries,text=black!75,anchor=north west},
  arr/.style={-{Stealth[length=1.6mm]},black!50,line width=0.55pt}]
  \fill[black!4,rounded corners=3pt] (-0.15,0.5) rectangle (4.85,-5.35);
  \fill[black!4,rounded corners=3pt] (5.6,0.5) rectangle (10.6,-5.35);
  \fill[black!4,rounded corners=3pt] (11.35,0.5) rectangle (16.35,-5.35);
  \node[ph] at (-0.05,0.42) {Phase 1: shared artefacts};
  \node[ph] at (5.7,0.42) {Phase 2: studies using the traces};
  \node[ph] at (11.45,0.42) {Phase 3: traces and QoE model};
  \node[pj,fill=gapcolor!7,draw=gapcolor!70!black] (a1) at (2.35,-0.5) {\textbf{A1} XR-Trace: open, frame-synchronised traces checked against 3GPP models};
  \node[pj,fill=gapcolor!7,draw=gapcolor!70!black] (a2) at (2.35,-1.75) {\textbf{A2} Rated trace-replay dataset for interactive XR QoE};
  \node[pj] (a4) at (2.35,-4.6) {\textbf{A4} Measured GPU and headset-energy cost models for split rendering};
  \foreach \n/\y/\t in {a3/-0.45/{\textbf{A3} Closed-loop benchmark for viewport and pose predictors},
                        a5/-1.32/{\textbf{A5} L4S, delay-based control and FEC on real bottlenecks},
                        a6/-2.19/{\textbf{A6} Privacy--QoE frontier of traffic defences},
                        a7/-3.06/{\textbf{A7} Haptic tolerance curves and an open haptic trace set},
                        a8/-3.93/{\textbf{A8} Interest-managed avatar replication beyond 100 clients},
                        a9/-4.8/{\textbf{A9} Percentile latency budgets for 6DoF split rendering}}
    {\node[pj] (\n) at (8.1,\y) {\t};}
  \foreach \n/\y/\t in {a10/-0.45/{\textbf{A10} Uplink-first scheduling for dense co-located headsets},
                        a11/-1.5/{\textbf{A11} Standards-conformant immersive streaming over 5G},
                        a12/-2.55/{\textbf{A12} Sim-to-real transfer of learned controllers on an open RAN},
                        a13/-3.6/{\textbf{A13} Group QoE and fairness for three to six users},
                        a14/-4.65/{\textbf{A14} Semantic, splat and geometric telepresence at matched bitrate}}
    {\node[pj] (\n) at (13.85,\y) {\t};}
  \coordinate (b1) at (5.22,0);
  \draw[arr] (a1.east) -- (b1 |- a1.east) |- (a3.west);
  \foreach \t in {a5,a6,a7,a8,a9}{\draw[arr] (b1 |- a1.east) |- (\t.west);}
  \draw[arr] (a1.south) -- (a2.north);
  \coordinate (b2) at (10.97,0);
  \draw[arr,gapcolor!70!black,rounded corners=2pt] (a2.east) -- ++(0.25,0) -- ++(0,-3.85) -- (10.97,-5.6+0.0) -- (b2 |- a14.west);
  \foreach \t in {a10,a11,a12,a13,a14}{\draw[arr,gapcolor!70!black] (10.97,-5.1) |- (\t.west);}
\end{tikzpicture}\end{adjustbox}
\caption{Research agenda as a dependency roadmap. Phase-1 projects (red) build the artefacts that the
others consume. A1 supplies traces to every later project (grey arrows); A2, built on A1, adds the
validated QoE model that the Phase-3 projects optimise against (red arrows). A4 stands alone.}\label{fig:roadmap}
\end{figure}

\paragraph{Phase 1: build what others need.} \emph{XR-Trace} (A1) would instrument the main workloads on
5G standalone and Wi-Fi~6E/7, with one, four and eight users. These are split-rendered VR, AR camera
offloading, avatar and shared-state traffic, tiled 360-degree video and twin-state updates. Each trace
would record pose, encoder output, packets at both ends, application events and photodiode MTP on a
common clock. The project would then fit traffic models, compare them with the 3GPP XR models, and release
the traces with a replay harness for ns-3 and NetEm. It comes first because every
later project consumes it. The \emph{rated replay dataset} (A2) would have
participants perform interactive tasks while the network replays A1's traces. It would collect momentary
and session ratings, sickness scores, task performance and packet logs together, so that a QoE model can
be fitted and then used as the objective by other work. It would also measure the difference between
replayed dynamics and independent impairments of the same mean. \emph{Energy cost models} (A4) would
replace the assumed GPU and battery costs inside offloading optimisers with measured ones.

\paragraph{Phase 2 and 3: answer the open questions on shared ground.} With traces in hand, five
studies become cheap and comparable. They are a closed-loop predictor benchmark with persistence and
Kalman baselines (A3), transport on real bottlenecks (A5), a privacy--QoE frontier for padding and jitter
defences (A6), haptic tolerance curves (A7) and interest-managed replication beyond one hundred clients
(A8). A percentile latency budget for 6DoF split rendering (A9) completes the phase. With a validated QoE
model as well, the remaining projects can optimise against experience rather than a proxy. These are
uplink-first scheduling (A10), a standards-conformant 5G pipeline (A11), sim-to-real transfer of learned
controllers as xApps on an open RAN (A12), group QoE and fairness (A13), and a matched-bitrate comparison
of semantic, Gaussian-splat and geometric telepresence (A14).

\paragraph{The evaluation bar.} Each project should meet the standard that Transactions-level reviewers
apply, and that much current work does not yet meet. First, a \emph{baseline ladder}: a trivial baseline,
a classical method, the strongest recent method re-implemented under the same conditions, an ablation
and, where computable, an oracle. Second, \emph{realism with control}: real hardware for the claim and
trace replay for repeatability, with the gap between simulation and measurement reported. Third,
\emph{out-of-distribution tests} for anything learned, on held-out users, content and networks. Fourth,
\emph{sound statistics}: repeated runs with confidence intervals and, for user studies, a power analysis,
within-subjects designs and mixed-effects models. Fifth, \emph{reported cost and released artefacts}: run
time, energy and signalling overhead on the target device, with code, traces and analysis scripts
published.
